\chapter{System Requirements and Architecture}\label{ch:system-requirements-and-architecture}

\section{Requirements Development}\label{requirements-development}

Requirements translate the research questions into observable evidence.
The repository does not contain an approved numerical specification for
pitch error, speed, runtime, or thermal margin. Numerical thresholds in
Table 3.1 are therefore identified as author- or advisor-controlled
values rather than silently selected targets. A requirement becomes
verified only when its method, instrumentation, operating condition, and
result are recorded.

\clearpage
\begin{landscape}
\par\medskip\noindent\singlespacing
\phantomsection\addcontentsline{lot}{table}{\protect\numberline{3.1}System requirements and verification matrix}
\textbf{Table 3.1. System requirements and verification matrix}\par

\begin{longtable}[]{@{}
  >{\raggedright\arraybackslash}p{(\columnwidth - 8\tabcolsep) * \real{0.0846}}
  >{\raggedright\arraybackslash}p{(\columnwidth - 8\tabcolsep) * \real{0.2615}}
  >{\raggedright\arraybackslash}p{(\columnwidth - 8\tabcolsep) * \real{0.2538}}
  >{\raggedright\arraybackslash}p{(\columnwidth - 8\tabcolsep) * \real{0.2077}}
  >{\raggedright\arraybackslash}p{(\columnwidth - 8\tabcolsep) * \real{0.1923}}@{}}
\toprule\noalign{}
\begin{minipage}[b]{\linewidth}\raggedright
\textbf{ID}
\end{minipage} & \begin{minipage}[b]{\linewidth}\raggedright
\textbf{Requirement}
\end{minipage} & \begin{minipage}[b]{\linewidth}\raggedright
\textbf{Acceptance criterion}
\end{minipage} & \begin{minipage}[b]{\linewidth}\raggedright
\textbf{Verification method}
\end{minipage} & \begin{minipage}[b]{\linewidth}\raggedright
\textbf{Current status}
\end{minipage} \\
\midrule\noalign{}
\endhead
\bottomrule\noalign{}
\endlastfoot
FR-01 & The robot shall command four independently addressable
actuators. & All four installed node IDs respond and report valid state.
& CAN enumeration and command\slash{}feedback test. & Partial: four slots
exist; only the verified active IDs may be claimed. \\
FR-02 & The system shall estimate body pitch and pitch rate. & Correct
sign and scale through the required angle\slash{}rate range; freshness
monitored. & Static orientations, controlled rotations, and timestamp
audit. & Partial: acquisition and estimator infrastructure exist; final
calibration pending. \\
FR-03 & The completed controller shall maintain upright balance. &
\textless AUTHOR INPUT REQUIRED: RMS pitch, peak excursion, settling
time, and test duration limits\textgreater. & Restrained then
free-standing repeat trials. & Not demonstrated in available
evidence. \\
FR-04 & The wheel arrangement shall provide commanded x, y, and yaw
motion while upright. & Measured geometry matrix rank equals 3 and
commanded signs match measured motion. & Analytical rank test,
wheel-direction test, velocity trials. & Pending geometry and balance
validation. \\
PR-01 & The real-time loop shall meet its deadline. & \textless AUTHOR
INPUT REQUIRED: period, maximum execution time, and jitter
limit\textgreater. & Timestamp or logic-analyzer measurement under
worst-case load. & Not measured. \\
PR-02 & CAN communication shall remain within the allowed age and error
limits. & \textless AUTHOR INPUT REQUIRED: feedback age, timeout,
bus-off, and frame-loss limits\textgreater. & Bus logging and fault
injection. & Basic motor communication demonstrated; statistics
absent. \\
ER-01 & Every device shall operate within its rated voltage range. &
Measured rail min\slash{}max within component limits during startup, idle, and
peak load. & DMM and oscilloscope measurements at load. & Design review
in progress; final PCB unverified. \\
ER-02 & Power conductors and protection shall support the selected
actuator limits. & Temperature rise and voltage drop remain below
approved limits; faults clear safely. & Current-limited test, load test,
thermal imaging, fuse test. & PCB and protection implementation
pending. \\
SR-01 & Loss of valid command, IMU state, or actuator feedback shall
remove torque permission. & Command becomes zero\slash{}disabled within
\textless AUTHOR INPUT REQUIRED: timeout\textgreater. & Unplug,
stale-data, and injected-fault tests. & Safety-gated architecture
reported; final validation pending. \\
MR-01 & Firmware shall separate portable logic from platform-specific
I\slash{}O. & Common modules compile independently and platform code owns
hardware access. & Repository structure and build review. & Implemented
in current architecture at a structural level. \\
\end{longtable}
\end{landscape}
\clearpage

\emph{Note.} FR = functional, PR = performance, ER = electrical, SR =
safety, and MR = maintainability requirement.

\section{System Architecture}\label{system-architecture}

The system is divided into mechanical, power, sensing, computation,
communication, user-interface, and safety subsystems. The mechanical
structure carries the battery, electronics, and body mass above the
common wheel axis. Four independent actuators connect the controller to
the ground through the Mecanum wheels. The IMU measures body motion; the
radio receiver provides operator commands; the display provides local
status; and CAN exchanges commands and feedback with the actuators.

Power enters from the six-cell battery. The actuator branches use the
battery domain directly, subject to fusing, switching, and protection
that must be confirmed in the final schematic. A buck converter produces
the logic rail, and local regulation or the controller-board regulator
produces 3.3 V where required. Logic ground and power return share a
controlled reference at the system level, but high-current motor return
paths must not force switching current through sensitive sensor or
transceiver return paths.

\begin{longtable}[]{@{}
  >{\raggedright\arraybackslash}p{(\columnwidth - 0\tabcolsep) * \real{1.0000}}@{}}
\toprule\noalign{}
\begin{minipage}[b]{\linewidth}\raggedright
\textless INSERT FIGURE 3.1: Draw a system block diagram containing the
6S battery, master disconnect and emergency stop, branch fuses, optional
BMS\slash{}protection board, actuator power branches, logic buck converter,
ESP32 prototype and STM32 target, ISM330DHCX IMU, CAN transceiver, four
AK40-10 actuators, RadioMaster receiver, ST7789 display, and
debug\slash{}logging interface\textgreater{}\\
What the figure should show: Draw a system block diagram containing the
6S battery, master disconnect and emergency stop, branch fuses, optional
BMS\slash{}protection board, actuator power branches, logic buck converter,
ESP32 prototype and STM32 target, ISM330DHCX IMU, CAN transceiver, four
AK40-10 actuators, RadioMaster receiver, ST7789 display, and
debug\slash{}logging interface\\
Likely source: KiCad schematic, firmware architecture, and final
hardware inventory\strut
\end{minipage} \\
\midrule\noalign{}
\endhead
\bottomrule\noalign{}
\endlastfoot
\end{longtable}

\begin{center}\singlespacing
\phantomsection\addcontentsline{lof}{figure}{\protect\numberline{3.1}System-level architecture of the Zephyr robot.}
\textbf{Figure 3.1. System-level architecture of the Zephyr robot.}
\end{center}

Information flows in the opposite direction from control authority.
Drivers convert peripheral data into timestamped measurements.
Estimation converts calibrated IMU data into pitch and rate. A
supervisor decides whether the current operating mode permits torque.
The controller produces a desired body wrench or wheel command. An
allocation stage maps the body command to four actuators. The motor task
transmits only commands that have passed safety checks. Feedback and
faults return through CAN and update the supervisor. This ownership
prevents the radio, UI, or a bring-up test from bypassing the safety
decision in production operation.

\begin{longtable}[]{@{}
  >{\raggedright\arraybackslash}p{(\columnwidth - 0\tabcolsep) * \real{1.0000}}@{}}
\toprule\noalign{}
\begin{minipage}[b]{\linewidth}\raggedright
\textless INSERT FIGURE 3.2: Use separate arrow styles for battery
power, regulated power, raw sensor data, estimated state, operator
command, safety permission, controller output, CAN command, CAN
feedback, faults, and telemetry\textgreater{}\\
What the figure should show: Use separate arrow styles for battery
power, regulated power, raw sensor data, estimated state, operator
command, safety permission, controller output, CAN command, CAN
feedback, faults, and telemetry\\
Likely source: firmware module diagram and electrical interface
diagram\strut
\end{minipage} \\
\midrule\noalign{}
\endhead
\bottomrule\noalign{}
\endlastfoot
\end{longtable}

\begin{center}\singlespacing
\phantomsection\addcontentsline{lof}{figure}{\protect\numberline{3.2}Power, sensing, command, and feedback paths among major subsystems.}
\textbf{Figure 3.2. Power, sensing, command, and feedback paths among major subsystems.}
\end{center}

\section{Subsystem Interfaces}\label{subsystem-interfaces}

\clearpage
\begin{landscape}
\par\medskip\noindent\singlespacing
\phantomsection\addcontentsline{lot}{table}{\protect\numberline{3.2}Major subsystem interfaces}
\textbf{Table 3.2. Major subsystem interfaces}\par

\begin{longtable}[]{@{}
  >{\raggedright\arraybackslash}p{(\columnwidth - 8\tabcolsep) * \real{0.1615}}
  >{\raggedright\arraybackslash}p{(\columnwidth - 8\tabcolsep) * \real{0.1615}}
  >{\raggedright\arraybackslash}p{(\columnwidth - 8\tabcolsep) * \real{0.1615}}
  >{\raggedright\arraybackslash}p{(\columnwidth - 8\tabcolsep) * \real{0.2077}}
  >{\raggedright\arraybackslash}p{(\columnwidth - 8\tabcolsep) * \real{0.3077}}@{}}
\toprule\noalign{}
\begin{minipage}[b]{\linewidth}\raggedright
\textbf{Producer}
\end{minipage} & \begin{minipage}[b]{\linewidth}\raggedright
\textbf{Consumer}
\end{minipage} & \begin{minipage}[b]{\linewidth}\raggedright
\textbf{Interface}
\end{minipage} & \begin{minipage}[b]{\linewidth}\raggedright
\textbf{Information or power}
\end{minipage} & \begin{minipage}[b]{\linewidth}\raggedright
\textbf{Required verification}
\end{minipage} \\
\midrule\noalign{}
\endhead
\bottomrule\noalign{}
\endlastfoot
6S battery\slash{}protection & Four actuator branches & VBAT and power return &
High-current energy & Polarity, fuse rating, conductor loss, peak droop,
disconnect behavior. \\
Logic converter & Controller and peripherals & 5 V and 3.3 V rails &
Regulated power & Startup, ripple, load transient, thermal rise,
undervoltage. \\
ISM330DHCX & MCU & SPI plus interrupt lines & Acceleration, angular
rate, data-ready & WHO\_AM\_I, rate, range, axes, timestamp, missed
samples. \\
RadioMaster RP2 V2 & MCU & CRSF UART & Operator channels and link state
& Baud rate, TX\slash{}RX crossing, failsafe, frame timeout. \\
MCU & ST7789 display & SPI and control GPIO & Status and diagnostics &
Voltage, pin mapping, refresh load, no control-loop blocking. \\
MCU\slash{}TWAI & TJA1051T\slash{}3 & TXD\slash{}RXD & CAN logic & 3.3 V compatibility,
standby state, error behavior. \\
CAN transceiver & AK40-10 network & CANH\slash{}CANL and reference & Commands,
feedback, faults & Linear topology, two terminations, bitrate, IDs,
timeouts. \\
Estimator & Supervisor\slash{}control & Timestamped snapshot & Pitch, rates,
validity and age & Atomic access, freshness, calibration status. \\
Supervisor\slash{}safety & Motor command path & Mode and torque permission &
Authorization and zero-torque request & Priority, deterministic fault
response, no bypass. \\
\end{longtable}
\end{landscape}
\clearpage

\section{Design Tradeoffs}\label{design-tradeoffs}

Four actuators provide redundancy and independent force allocation, but
they increase current demand, CAN traffic, wiring, configuration effort,
and the number of possible sign errors. Mecanum wheels provide lateral
authority, but their rollers introduce compliance, vibration, changing
contact conditions, and slip. Integrated actuators reduce the burden on
the central MCU, but their internal loop settings and motor-side encoder
cannot be treated as an ideal wheel-torque source. A custom PCB can
improve wiring and grounding, but it creates manufacturing risk if the
schematic is released before the firmware pin contract and power-tree
design are frozen.

The ESP32 is suitable for rapid bring-up because the current firmware
and peripheral interfaces are already developed on that platform. The
STM32H723ZG offers greater deterministic-control capability and FDCAN
resources \cite{st2026stm32h723}, but migration creates a second implementation that
must reproduce the same common logic and safety behavior. The selected
strategy is to stabilize architecture and interfaces in portable
modules, complete hardware evidence on one platform, and migrate only
when the control and timing requirements justify the change.

\section{Architecture Status}\label{architecture-status}

The repository evidence supports a modular architecture and several
functioning device paths, but not a completed system. The mechanical
assembly is reported to be approximately 80 percent complete and the
motors have been moved from the prototype firmware. The electrical
design has progressed from point-to-point wiring toward a custom board,
but the reviewed PCB state was unrouted or pre-fabrication. The firmware
contains sensing, communication, UI, supervisor, and safety
infrastructure, but the final balance law and wheel allocation remain
absent or unverified. The simulation work provides a symbolic and
numerical starting point, not experimentally validated results. These
status distinctions govern the wording of the remaining chapters.
